\scanpage{82}
\mspage{75}
% archive sheet number: 81
et de plus \inserted{\unsure{avec}} $C,\mathrm{TC}$ \deleted{\unsure{de type fini}} tels qu'on ait
\[
C=\widetilde{C}_0,\qquad \mathrm{TC}=\widetilde{\mathrm{TC}}_0,
\]
avec $C_0,\mathrm{TC}_0$ petits. On pose maintenant
% slanted
\margin{On dit des $C,\mathrm{TC}$ (\unsure{en tant} que parties \unsure{closes} à droite) ``\unsure{accessibles|coaccessibles}''.}
\[
\begin{aligned}
F(I)&=\{f\in\mathrm{Fl}\,M(I)\mid f_s\in F\quad \forall s\in\mathrm{Ob}\,I\}\,,\\
\mathrm{TF}(I)&=\{f\in\mathrm{Fl}\,M(I)\mid f_s\in\mathrm{TF}\quad \forall s\in\mathrm{Ob}\,I\}\,,
\end{aligned}
\]
puis
\[
\begin{aligned}
C(I)&=\mathrm{TF}(I)^{\unsure{\alpha|\ast}},\\
\mathrm{TC}(I)&=F(I)^\alpha.
\end{aligned}
\]
L'espoir (que je devrai vérifier ultérieurement), c'est qu'on obtient ainsi une quillénisation uniforme à droite \inserted{de $W$,} sur $\mathrm{Cat}_{\unsure{\mathcal U|\mathscr U}}$ \unsure{tout entier}.

Le seul problème, c'est la validité des axiomes de factorisation de $f\in\mathrm{Fl}\,M(I)$ en $f=\unsure{pi|pj}$, avec $i\in\mathrm{TC}$, $\unsure{f|p}\in F$, et aussi avec $i\in C$, $f\in\mathrm{TF}$. J'ai bon espoir que c'est vrai --- cela se réduit essentiellement à voir pour $(C(I),\mathrm{TF}(I))$ \deleted{\unsure{soit}} et $(\mathrm{TC}(I),F(I))$, \deleted{\unsure{on doit le}} soit $(\unsure{\Phi|\varphi}(I),\unsure{\Psi|\Upsilon}(I))$, \deleted{\unsure{on doit prouver|on suppose}} en plus de $\Phi(I)=\Psi(I)^\alpha$, à $\Psi(I)=\Phi(I)_{\unsure{\alpha\delta|\alpha\beta}}$ (i.e. $\Psi(I)$ est clos à droite), et que $\Phi(I)$ soit accessible (tout comme $\Phi=\Phi(e)$).

Si ça marche, on aura l'\deleted{\unsure{admissibilité}}
